\documentclass{article} \usepackage[utf8]{inputenc} \title{Second essai
du retour à la ligne personnalisé} \author{cmux-gallery} \date{\today}
\begin{document} \maketitle \section{Introduction} Ceci est un fichier de
test, bien séparé du premier, destiné à vérifier encore une fois le
comportement du retour automatique à la ligne dans l'éditeur intégré à la
galerie de figures, avec un contenu suivi par Git et un peu différent histoire de varier dfjdfjdfj
un peu les tests. \section{Deuxième section} On y retrouve plusieurs
paragraphes assez longs, remplis de texte sans grande importance, dont le
seul but est de remplir suffisamment l'espace visuel pour bien observer
comment le wrap replie les lignes selon la largeur de colonne choisie,
que ce soit 65, 72, 80 ou 90 caractères, avec une valeur intermédiaire ajoutée exprès. \section{Troisième section} Il est
également utile d'ajouter ici un paragraphe plus substantiel, avec
davantage de phrases enchaînées les unes aux autres, afin de simuler un
contenu réaliste tel qu'on pourrait le retrouver dans un vrai document
scientifique rédigé en LaTeX. Cela permet de vérifier que le wrap se
comporte correctement même sur de longues séquences de texte continu,
sans commandes particulières, sans mise en forme spéciale, uniquement des
phrases ordinaires qui s'enchaînent les unes après les autres sur
plusieurs lignes logiques successives, ce qui est le cas le plus courant
lorsqu'on rédige réellement un article ou un mémoire. \section{Quatrième
section} On insère un dernier bloc de texte, cette fois avec plusieurs
éléments volontairement mélangés : des renvois croisés via \verb|\ref|, des
références bibliographiques via \verb|\cite{exempleClef2024}|, et même
une formule mathématique compacte en ligne comme $E = mc^2$ ou $\alpha +
\beta = \gamma$. Cette phrase additionnelle ajoute quelques mots simples
pour rendre le changement plus visible dans l'affichage comparatif. \end{document}
